\chapter{KV Cache 性能基准测试与吞吐加速比实测}

\section{Lab 06：KV Cache 原理与自回归解码基准测试}

\section{[思考] 为什么需要 KV Cache？}

在大语言模型（如 GPT、Llama、Qwen 等 Decoder-only 架构）生成文本时，是\textbf{逐字（Token-by-Token）生成}的：

\begin{lstlisting}[numbers=none]
Prompt: "人工智能的未来是"
Step 1 -> "星"
Step 2 -> "辰"
Step 3 -> "大"
Step 4 -> "海"
\end{lstlisting}

在标准的 Self-Attention 计算中：

\[
\text{Attention}(Q, K, V) = \text{softmax}\left(\frac{Q K^T}{\sqrt{d_k}}\right) V
\]

\subsection{朴素实现的问题（无 KV Cache）：}

生成第 $t$ 个 token 时，把前 $t-1$ 个 token 连同最新 token 一起重新送入模型：
\begin{itemize}
  \item 重新计算前 $t-1$ 个 token 的 $Q, K, V$ 投影矩阵。
  \item 但事实上，\textbf{历史 token 的 Key 和 Value 向量是固定不变的}！
  \item 每次生成都会重复计算历史所有 token 的 $K$ 和 $V$，导致总计算复杂度高达 $O(N^2)$，生成越到后面越慢。
\end{itemize}

\subsection{KV Cache 的优化思路：}

\begin{itemize}
  \item 在 \textbf{Prefill 阶段}（首字），计算出 Prompt 中所有 token 的 $K$ 和 $V$，并保存在显存中（KV Cache）。
  \item 在 \textbf{Decode 阶段}（后续生成），每次\textbf{只输入最新的 1 个 token}：
  \item 只对最新 token 计算 $q_{\text{new}}, k_{\text{new}}, v_{\text{new}}$。
  \item 把 $k_{\text{new}}, v_{\text{new}}$ 拼接到缓存列表中：$K_{\text{all}} = [K_{\text{cache}}, k_{\text{new}}]$。
  \item 用单个查询向量 $q_{\text{new}}$ 与整个 $K_{\text{all}}$ 做点积得到注意力分数，加权聚合 $V_{\text{all}}$。
  \item 这样，单步生成的计算复杂度从 $O(t)$ 矩阵乘降为向量-矩阵乘，避免了大量冗余算力消耗！
\end{itemize}

\vspace{0.8em}\hrule\vspace{0.8em}

\section{[执行] 运行测试}

本目录下提供了 \textbf{\texttt{kv\_cache\_benchmark.py}}，演示：
\begin{enumerate}
  \item 朴素生成 vs KV Cache 生成的数学等价性（验证输出完全一致）。
  \item 每步生成耗时随序列长度增加的变化趋势对比。
\end{enumerate}

\begin{lstlisting}[language=bash]
python3 lab06_kv_cache/kv_cache_benchmark.py
\end{lstlisting}

\textit{(本脚本自适应纯 Python 或 NumPy，即使暂未安装任何第三方库也能直接运行！)}

\vspace{0.8em}\hrule\vspace{0.8em}

\section{[链接] 继续深入}

本模块用的是随机权重的单层注意力。想在\textbf{真正训练过的模型}上验证同样的结论（输出逐 token 一致、缓存字节数与公式逐字节相等、每步耗时对比），请做：
\begin{itemize}
  \item {}\textbf{Lab 05}：自己训练的迷你 Llama（\texttt{generate.py}）；
  \item {}\textbf{Lab 07b}：真实的 Qwen2.5-0.5B，上下文越长 KV Cache 的优势越大（2048 时约快一个数量级）。
\end{itemize}

注意：KV Cache 消除的是对历史 token 重复做投影 / FFN 的冗余，\textbf{注意力本身每步仍要读完整个缓存}，长上下文时这会成为新的瓶颈（第 6 篇 §3 的补充说明、第 8 篇 §4.3）。


\section{实验核心源码精选与解析}

本实验配套有完整可运行的工业级验证代码，重点实现细节如下：


\subsection{源码实现：\texttt{kv\_cache\_benchmark.py}}

\lstinputlisting[language=Python, caption={\texttt{kv\_cache\_benchmark.py}}]{code/lab06_kv_cache_kv_cache_benchmark.py}

